\begin{abstract}
Clinicians need to know which human-interpretable
quantitative axes its representations encode, whether those axes survive cohort and technical
variation, and whether downstream decisions use them. We audited six non-interchangeable axes
in frozen CONCH and Virchow: Grade/ISUP, tumor phenotype/content, PTEN loss, AR
activity, SPOP mutation, and recurrence. Among the six axes, grade and phenotype showed the strongest transport;
PTEN and AR were conditional, SPOP was unsupported, and recurrence changed
with endpoint and reference. ISUP alone underwent locked-head functional tests. Internal targeted
erasure reduced two biochemical-recurrence heads' concordance beyond matched-random controls. In
site-heldout TCGA, only Virchow passed. In independent-patient LEOPARD, neither encoder qualified.
In CHIMERA (95 patients; 27 events), ISUP was recoverable in both encoders,
but only Virchow met the prespecified external whole-tissue functional-transport gate; CONCH
erasure was positive while its full-head interval reached chance. Functional transport was
therefore model-specific. The comparative map provides an actionable
audit vocabulary: transported coordinates can anchor clinician review of agreement and
disagreement, conditional or unsupported axes indicate where explanations should be withheld,
and gaps prioritize external or functional validation. It enables
clinician-understandable reliability audits and bounds future residual-signal discovery; it does not
establish clinical benefit, tumor-specific mechanisms, encoder superiority, or new biomarkers.
\end{abstract}
